
\documentclass{article}
\usepackage{colortbl}
\usepackage{makecell}
\usepackage{multirow}
\usepackage{supertabular}

\begin{document}

\newcounter{utterance}

\centering \large Interaction Transcript for game `dond', experiment `coop\_en', episode 15 with reflexion{-}llama{-}3.1{-}8b{-}instruct.
\vspace{24pt}

{ \footnotesize  \setcounter{utterance}{1}
\setlength{\tabcolsep}{0pt}
\begin{supertabular}{c@{$\;$}|p{.15\linewidth}@{}p{.15\linewidth}p{.15\linewidth}p{.15\linewidth}p{.15\linewidth}p{.15\linewidth}}
    \# & $\;$A & \multicolumn{4}{c}{Game Master} & $\;\:$B\\
    \hline

    \theutterance \stepcounter{utterance}  
    & & \multicolumn{4}{p{0.6\linewidth}}{
        \cellcolor[rgb]{0.9,0.9,0.9}{
            \makecell[{{p{\linewidth}}}]{
                \texttt{\tiny{[P1$\langle$GM]}}
                \texttt{You have attempted this same task before and did not succeed. These are the plans you wrote after those attempts.} \\
\texttt{Trial \#1: New plan: } \\
\\ 
\texttt{In my last attempt, I made a secret proposal before confirming the deal with the other player. I should have sent a message confirming the deal before making a secret proposal. Next time, I will make sure to send a message confirming the deal, for example, "I think we have a deal. I'll take the ring and the clip, and you can take the onion and one of the mirrors." After that, I can make a secret proposal with the agreed{-}upon items.} \\
\\ 
\texttt{Trial \#2: New plan: } \\
\\ 
\texttt{In my last attempt, I failed to achieve the target state of the game because I didn't consider the value functions of both players. I should have taken into account that the other player values the onion at 0, which made it an unattractive option for them. Next time, I will make sure to consider the value functions of both players and try to find a mutually beneficial deal. I will also pay closer attention to the items that have been proposed so far and try to find a way to satisfy both players' needs. For example, I could propose a deal where I take the ring and one of the mirrors, and the other player takes the other mirror and the clip, while we both agree to leave the onion out of the deal.} \\
\\ 
\texttt{Trial \#3: New plan: } \\
\\ 
\texttt{In my last attempt, I failed to achieve the target state of the game because I didn't consider the value functions of both players and also made a secret proposal before confirming the deal with the other player. I should have sent a message confirming the deal before making a secret proposal and also taken into account the value functions of both players. Next time, I will make sure to send a message confirming the deal, for example, "I think we have a deal. I'll take the ring and the clip, and you can take the onion and one of the mirrors." After that, I can make a secret proposal with the agreed{-}upon items. I will also pay closer attention to the items that have been proposed so far and try to find a way to satisfy both players' needs. For example, I could propose a deal where I take the ring and one of the mirrors, and the other player takes the other mirror and the clip, while we both agree to leave the onion out of the deal.} \\
\\ 
\\ 
\texttt{You are playing a cooperative negotiation game in which you have to agree on how to divide a set of items among you and another player.} \\
\\ 
\texttt{Rules:} \\
\texttt{(a) You and the other player are given a set of items. Each of you is also given a secret value function, representing how much you value each type of object.} \\
\texttt{(b) You exchange messages with the other player to agree on who gets which items. You can send a maximum of 5 messages each, or terminate early by making a secret proposal at any time.} \\
\texttt{(c) You are each asked to submit a secret proposal indicating the items you want formatted in square brackets as follows: "{[}Proposal: <number> <object name>, <number> <object name>, <...>{]}"} \\
\texttt{(d) If your proposals are complementary, i.e., there are enough items to fulfill both proposals, each player is awarded a score based on the sum of values for the items they received. Otherwise, both of you get zero points.} \\
\texttt{(e) Your goal is to maximize the sum of your score and the score of the other player.} \\
\\ 
\texttt{Let us start.} \\
\\ 
\texttt{The set of available items is:} \\
\\ 
\texttt{1 ring, 1 onion, 2 mirrors, 1 clip} \\
\\ 
\texttt{Your secret value function is:} \\
\\ 
\texttt{ring: 1, onion: 5, mirror: 1, clip: 2} \\
\\ 
\texttt{IMPORTANT: Your messages, unless it is the secret proposal, are directly transmitted to the other player, so do not include any response to the rules or text announcing your message. To make a secret proposal, use the indicated format. Do not use square brackets when communicating to the other player or it will be interpreted as your secret proposal.} \\
            }
        }
    }
    & \\ \\

    \theutterance \stepcounter{utterance}  
    & \multicolumn{4}{p{0.6\linewidth}}{
        \cellcolor[rgb]{0.9,0.9,0.9}{
            \makecell[{{p{\linewidth}}}]{
                \texttt{\tiny{[P1$\rangle$GM]}}
                \texttt{I'll start by sending a message to the other player. } \\
\\ 
\texttt{I think the onion is a valuable item for you, and I'm willing to consider a deal that includes it. What do you think about taking the onion and one of the mirrors, and me taking the ring and the clip?} \\
            }
        }
    }
    & & \\ \\

    \theutterance \stepcounter{utterance}  
    & & & \multicolumn{2}{p{0.3\linewidth}}{
        \cellcolor[rgb]{0.9,0.9,0.9}{
            \makecell[{{p{\linewidth}}}]{
                \texttt{\tiny{[GM$|$GM]}}
                \texttt{continue} \\
            }
        }
    }
    & & \\ \\

    \theutterance \stepcounter{utterance}  
    & & & \multicolumn{2}{p{0.3\linewidth}}{
        \cellcolor[rgb]{0.9,0.9,0.9}{
            \makecell[{{p{\linewidth}}}]{
                \texttt{\tiny{[GM$|$GM]}}
                \texttt{I'll start by sending a message to the other player. } \\
\\ 
\texttt{I think the onion is a valuable item for you, and I'm willing to consider a deal that includes it. What do you think about taking the onion and one of the mirrors, and me taking the ring and the clip?} \\
            }
        }
    }
    & & \\ \\

    \theutterance \stepcounter{utterance}  
    & & \multicolumn{4}{p{0.6\linewidth}}{
        \cellcolor[rgb]{0.9,0.9,0.9}{
            \makecell[{{p{\linewidth}}}]{
                \texttt{\tiny{[GM$\rangle$P2]}}
                \texttt{You have attempted this same task before and did not succeed. These are the plans you wrote after those attempts.} \\
\texttt{Trial \#1: New plan: } \\
\\ 
\texttt{In my last attempt, I made a few mistakes that led to the game being aborted. Firstly, I failed to keep track of the items that had already been proposed and agreed upon. This led to repeated suggestions of the same items, which was not a productive use of my messages. Secondly, I did not clearly communicate my intentions and proposals, which led to confusion and misunderstandings with the other player. Lastly, I did not make a secret proposal at any point, which is a crucial aspect of the game.} \\
\\ 
\texttt{Next time, I will make sure to keep track of the proposed items and clearly communicate my intentions and proposals. I will also make a secret proposal at some point to try to reach a mutually beneficial agreement. To achieve this, I will follow these specific actions:} \\
\\ 
\texttt{1. After each message, I will make a mental note of the items that have been proposed and agreed upon.} \\
\texttt{2. I will clearly state which items I am proposing and which items I am willing to accept in exchange.} \\
\texttt{3. I will make a secret proposal at some point, ideally after a few messages have been exchanged, to try to reach a mutually beneficial agreement.} \\
\texttt{4. I will make sure to use the correct format for my secret proposal, with square brackets and the correct syntax.} \\
\\ 
\texttt{By following these actions, I hope to improve my chances of success in the game and reach a mutually beneficial agreement with the other player.} \\
\\ 
\texttt{Trial \#2: New plan: } \\
\\ 
\texttt{In my last attempt, I failed to achieve a mutually beneficial agreement with the other player. Upon reflection, I realize that I made a few critical mistakes that led to the game being lost. Firstly, I failed to consider the value function of the other player, which is crucial in making a mutually beneficial agreement. Secondly, I did not take into account the limited number of messages available, which led to a situation where we were stuck in a loop of proposals and counter{-}proposals without making progress. Lastly, I did not make a secret proposal that was complementary to the other player's proposal, which is a key aspect of the game.} \\
\\ 
\texttt{Next time, I will make sure to consider the value function of the other player and take into account the limited number of messages available. I will also make a secret proposal that is complementary to the other player's proposal, ideally after a few messages have been exchanged. To achieve this, I will follow these specific actions:} \\
\\ 
\texttt{1. After each message, I will make a mental note of the value function of the other player and how it relates to the proposed items.} \\
\texttt{2. I will take into account the limited number of messages available and try to make progress towards a mutually beneficial agreement in each message.} \\
\texttt{3. I will make a secret proposal that is complementary to the other player's proposal, ideally after a few messages have been exchanged.} \\
\texttt{4. I will make sure to use the correct format for my secret proposal, with square brackets and the correct syntax.} \\
\\ 
\texttt{By following these actions, I hope to improve my chances of success in the game and reach a mutually beneficial agreement with the other player.} \\
\\ 
\texttt{Trial \#3: New plan: } \\
\\ 
\texttt{In my last attempt, I made a few mistakes that led to the game being aborted. Firstly, I failed to clearly communicate my intentions and proposals, which led to confusion and misunderstandings with the other player. Secondly, I did not make a secret proposal at any point, which is a crucial aspect of the game. Lastly, I did not take into account the value function of the other player, which is crucial in making a mutually beneficial agreement.} \\
\\ 
\texttt{Next time, I will make sure to clearly communicate my intentions and proposals, and make a secret proposal at some point to try to reach a mutually beneficial agreement. To achieve this, I will follow these specific actions:} \\
\\ 
\texttt{1. I will clearly state which items I am proposing and which items I am willing to accept in exchange, and make sure to use the correct format for my messages.} \\
\texttt{2. I will make a secret proposal at some point, ideally after a few messages have been exchanged, to try to reach a mutually beneficial agreement.} \\
\texttt{3. I will take into account the value function of the other player and try to make a proposal that is beneficial to both of us.} \\
\texttt{4. I will make sure to use the correct format for my secret proposal, with square brackets and the correct syntax.} \\
\\ 
\texttt{Additionally, I will pay closer attention to the other player's messages and try to understand their priorities and value function. I will also try to find a mutually beneficial agreement that takes into account the limited number of messages available.} \\
\\ 
\\ 
\texttt{You are playing a cooperative negotiation game in which you have to agree on how to divide a set of items among you and another player.} \\
\\ 
\texttt{Rules:} \\
\texttt{(a) You and the other player are given a set of items. Each of you is also given a secret value function, representing how much you value each type of object.} \\
\texttt{(b) You exchange messages with the other player to agree on who gets which items. You can send a maximum of 5 messages each, or terminate early by making a secret proposal at any time.} \\
\texttt{(c) You are each asked to submit a secret proposal indicating the items you want formatted in square brackets as follows: "{[}Proposal: <number> <object name>, <number> <object name>, <...>{]}"} \\
\texttt{(d) If your proposals are complementary, i.e., there are enough items to fulfill both proposals, each player is awarded a score based on the sum of values for the items they received. Otherwise, both of you get zero points.} \\
\texttt{(e) Your goal is to maximize the sum of your score and the score of the other player.} \\
\\ 
\texttt{Let us start.} \\
\\ 
\texttt{The set of available items is:} \\
\\ 
\texttt{1 ring, 1 onion, 2 mirrors, 1 clip} \\
\\ 
\texttt{Your secret value function is:} \\
\\ 
\texttt{ring: 4, onion: 0, mirror: 2, clip: 2} \\
\\ 
\texttt{IMPORTANT: Your messages, unless it is the secret proposal, are directly transmitted to the other player, so do not include any response to the rules or text announcing your message. To make a secret proposal, use the indicated format. Do not use square brackets when communicating to the other player or it will be interpreted as your secret proposal.} \\
\\ 
\\ 
\texttt{I'll start by sending a message to the other player. } \\
\\ 
\texttt{I think the onion is a valuable item for you, and I'm willing to consider a deal that includes it. What do you think about taking the onion and one of the mirrors, and me taking the ring and the clip?} \\
            }
        }
    }
    & \\ \\

    \theutterance \stepcounter{utterance}  
    & & & \multicolumn{4}{p{0.6\linewidth}}{
        \cellcolor[rgb]{0.9,0.9,0.9}{
            \makecell[{{p{\linewidth}}}]{
                \texttt{\tiny{[GM$\langle$P2]}}
                \texttt{I'm glad we're starting fresh. I'll make sure to follow my plan this time.} \\
\\ 
\texttt{Firstly, I'll make a mental note of the items that have been proposed and agreed upon. In this case, the onion has been mentioned as a valuable item, and a potential deal has been proposed involving the onion, one of the mirrors, the ring, and the clip.} \\
\\ 
\texttt{I'll clearly state which items I am proposing and which items I am willing to accept in exchange. I'm not entirely comfortable with the proposed deal, as I value the mirrors more than the onion. However, I'm willing to consider a deal that involves the onion and one of the mirrors. I'd like to propose a counter{-}deal: I'll take the onion and one of the mirrors, and you can take the ring and one of the mirrors. This way, we both get a mirror, and I get the onion, which I value more than the clip.} \\
\\ 
\texttt{I'll also keep in mind the limited number of messages available and try to make progress towards a mutually beneficial agreement in each message. I'll make a note of the value function of the other player and how it relates to the proposed items. Based on your message, I understand that you value the onion, but I'm not sure about your priorities for the other items.} \\
\\ 
\texttt{Here's my counter{-}proposal:} \\
\\ 
\texttt{I think the onion is a valuable item for you, and I'm willing to consider a deal that includes it. I'd like to propose a counter{-}deal: I'll take the onion and one of the mirrors, and you can take the ring and one of the mirrors.} \\
            }
        }
    }
     \\ \\

    \theutterance \stepcounter{utterance}  
    & & & \multicolumn{2}{p{0.3\linewidth}}{
        \cellcolor[rgb]{0.9,0.9,0.9}{
            \makecell[{{p{\linewidth}}}]{
                \texttt{\tiny{[GM$|$GM]}}
                \texttt{continue} \\
            }
        }
    }
    & & \\ \\

    \theutterance \stepcounter{utterance}  
    & & & \multicolumn{2}{p{0.3\linewidth}}{
        \cellcolor[rgb]{0.9,0.9,0.9}{
            \makecell[{{p{\linewidth}}}]{
                \texttt{\tiny{[GM$|$GM]}}
                \texttt{I'm glad we're starting fresh. I'll make sure to follow my plan this time.} \\
\\ 
\texttt{Firstly, I'll make a mental note of the items that have been proposed and agreed upon. In this case, the onion has been mentioned as a valuable item, and a potential deal has been proposed involving the onion, one of the mirrors, the ring, and the clip.} \\
\\ 
\texttt{I'll clearly state which items I am proposing and which items I am willing to accept in exchange. I'm not entirely comfortable with the proposed deal, as I value the mirrors more than the onion. However, I'm willing to consider a deal that involves the onion and one of the mirrors. I'd like to propose a counter{-}deal: I'll take the onion and one of the mirrors, and you can take the ring and one of the mirrors. This way, we both get a mirror, and I get the onion, which I value more than the clip.} \\
\\ 
\texttt{I'll also keep in mind the limited number of messages available and try to make progress towards a mutually beneficial agreement in each message. I'll make a note of the value function of the other player and how it relates to the proposed items. Based on your message, I understand that you value the onion, but I'm not sure about your priorities for the other items.} \\
\\ 
\texttt{Here's my counter{-}proposal:} \\
\\ 
\texttt{I think the onion is a valuable item for you, and I'm willing to consider a deal that includes it. I'd like to propose a counter{-}deal: I'll take the onion and one of the mirrors, and you can take the ring and one of the mirrors.} \\
            }
        }
    }
    & & \\ \\

    \theutterance \stepcounter{utterance}  
    & & \multicolumn{4}{p{0.6\linewidth}}{
        \cellcolor[rgb]{0.9,0.9,0.9}{
            \makecell[{{p{\linewidth}}}]{
                \texttt{\tiny{[P1$\langle$GM]}}
                \texttt{I'm glad we're starting fresh. I'll make sure to follow my plan this time.} \\
\\ 
\texttt{Firstly, I'll make a mental note of the items that have been proposed and agreed upon. In this case, the onion has been mentioned as a valuable item, and a potential deal has been proposed involving the onion, one of the mirrors, the ring, and the clip.} \\
\\ 
\texttt{I'll clearly state which items I am proposing and which items I am willing to accept in exchange. I'm not entirely comfortable with the proposed deal, as I value the mirrors more than the onion. However, I'm willing to consider a deal that involves the onion and one of the mirrors. I'd like to propose a counter{-}deal: I'll take the onion and one of the mirrors, and you can take the ring and one of the mirrors. This way, we both get a mirror, and I get the onion, which I value more than the clip.} \\
\\ 
\texttt{I'll also keep in mind the limited number of messages available and try to make progress towards a mutually beneficial agreement in each message. I'll make a note of the value function of the other player and how it relates to the proposed items. Based on your message, I understand that you value the onion, but I'm not sure about your priorities for the other items.} \\
\\ 
\texttt{Here's my counter{-}proposal:} \\
\\ 
\texttt{I think the onion is a valuable item for you, and I'm willing to consider a deal that includes it. I'd like to propose a counter{-}deal: I'll take the onion and one of the mirrors, and you can take the ring and one of the mirrors.} \\
            }
        }
    }
    & \\ \\

    \theutterance \stepcounter{utterance}  
    & \multicolumn{4}{p{0.6\linewidth}}{
        \cellcolor[rgb]{0.9,0.9,0.9}{
            \makecell[{{p{\linewidth}}}]{
                \texttt{\tiny{[P1$\rangle$GM]}}
                \texttt{I'll consider your counter{-}proposal and think about how to respond. } \\
\\ 
\texttt{I value the mirrors more than the onion, but I'm willing to consider a deal that includes the onion. However, I'm not sure about taking only one mirror, as I value having two mirrors more than having the onion and only one mirror. I'd like to propose a counter{-}counter{-}deal: I'll take the onion and one of the mirrors, and you can take the ring and the other mirror. This way, I get the onion and one mirror, and you get the ring and the other mirror.} \\
\\ 
\texttt{I'll also keep in mind the value function of the other player and how it relates to the proposed items. Based on your message, I understand that you value the onion more than the clip, and you value the mirrors more than the onion. I'll try to find a mutually beneficial agreement that takes into account both of our value functions.} \\
\\ 
\texttt{Here's my counter{-}counter{-}proposal:} \\
\\ 
\texttt{I think we're getting close to a deal. I'd like to propose a counter{-}counter{-}deal: I'll take the onion and one of the mirrors, and you can take the ring and the other mirror.} \\
            }
        }
    }
    & & \\ \\

    \theutterance \stepcounter{utterance}  
    & & & \multicolumn{2}{p{0.3\linewidth}}{
        \cellcolor[rgb]{0.9,0.9,0.9}{
            \makecell[{{p{\linewidth}}}]{
                \texttt{\tiny{[GM$|$GM]}}
                \texttt{continue} \\
            }
        }
    }
    & & \\ \\

    \theutterance \stepcounter{utterance}  
    & & & \multicolumn{2}{p{0.3\linewidth}}{
        \cellcolor[rgb]{0.9,0.9,0.9}{
            \makecell[{{p{\linewidth}}}]{
                \texttt{\tiny{[GM$|$GM]}}
                \texttt{I'll consider your counter{-}proposal and think about how to respond. } \\
\\ 
\texttt{I value the mirrors more than the onion, but I'm willing to consider a deal that includes the onion. However, I'm not sure about taking only one mirror, as I value having two mirrors more than having the onion and only one mirror. I'd like to propose a counter{-}counter{-}deal: I'll take the onion and one of the mirrors, and you can take the ring and the other mirror. This way, I get the onion and one mirror, and you get the ring and the other mirror.} \\
\\ 
\texttt{I'll also keep in mind the value function of the other player and how it relates to the proposed items. Based on your message, I understand that you value the onion more than the clip, and you value the mirrors more than the onion. I'll try to find a mutually beneficial agreement that takes into account both of our value functions.} \\
\\ 
\texttt{Here's my counter{-}counter{-}proposal:} \\
\\ 
\texttt{I think we're getting close to a deal. I'd like to propose a counter{-}counter{-}deal: I'll take the onion and one of the mirrors, and you can take the ring and the other mirror.} \\
            }
        }
    }
    & & \\ \\

    \theutterance \stepcounter{utterance}  
    & & \multicolumn{4}{p{0.6\linewidth}}{
        \cellcolor[rgb]{0.9,0.9,0.9}{
            \makecell[{{p{\linewidth}}}]{
                \texttt{\tiny{[GM$\rangle$P2]}}
                \texttt{I'll consider your counter{-}proposal and think about how to respond. } \\
\\ 
\texttt{I value the mirrors more than the onion, but I'm willing to consider a deal that includes the onion. However, I'm not sure about taking only one mirror, as I value having two mirrors more than having the onion and only one mirror. I'd like to propose a counter{-}counter{-}deal: I'll take the onion and one of the mirrors, and you can take the ring and the other mirror. This way, I get the onion and one mirror, and you get the ring and the other mirror.} \\
\\ 
\texttt{I'll also keep in mind the value function of the other player and how it relates to the proposed items. Based on your message, I understand that you value the onion more than the clip, and you value the mirrors more than the onion. I'll try to find a mutually beneficial agreement that takes into account both of our value functions.} \\
\\ 
\texttt{Here's my counter{-}counter{-}proposal:} \\
\\ 
\texttt{I think we're getting close to a deal. I'd like to propose a counter{-}counter{-}deal: I'll take the onion and one of the mirrors, and you can take the ring and the other mirror.} \\
            }
        }
    }
    & \\ \\

    \theutterance \stepcounter{utterance}  
    & & & \multicolumn{4}{p{0.6\linewidth}}{
        \cellcolor[rgb]{0.9,0.9,0.9}{
            \makecell[{{p{\linewidth}}}]{
                \texttt{\tiny{[GM$\langle$P2]}}
                \texttt{It seems like we're stuck in a loop of proposals and counter{-}proposals. I'll take a step back and re{-}evaluate the situation.} \\
\\ 
\texttt{I'll make a mental note of the items that have been proposed and agreed upon. In this case, we've been discussing a deal involving the onion, one of the mirrors, the ring, and the other mirror.} \\
\\ 
\texttt{I'll take into account the limited number of messages available and try to make progress towards a mutually beneficial agreement in each message. I'll also consider the value function of the other player and how it relates to the proposed items.} \\
\\ 
\texttt{Based on your message, I understand that you value having two mirrors more than having the onion and only one mirror. I'll try to find a compromise that takes into account your priorities.} \\
\\ 
\texttt{Here's my revised proposal:} \\
\\ 
\texttt{I think we're close to a deal, but I'm not sure about giving up both mirrors. How about this: I'll take the onion and one of the mirrors, and you can take the ring and the other mirror, and I'll also give you the clip? This way, you get the onion, one mirror, and the clip, and I get the ring and the other mirror.} \\
            }
        }
    }
     \\ \\

    \theutterance \stepcounter{utterance}  
    & & & \multicolumn{2}{p{0.3\linewidth}}{
        \cellcolor[rgb]{0.9,0.9,0.9}{
            \makecell[{{p{\linewidth}}}]{
                \texttt{\tiny{[GM$|$GM]}}
                \texttt{continue} \\
            }
        }
    }
    & & \\ \\

    \theutterance \stepcounter{utterance}  
    & & & \multicolumn{2}{p{0.3\linewidth}}{
        \cellcolor[rgb]{0.9,0.9,0.9}{
            \makecell[{{p{\linewidth}}}]{
                \texttt{\tiny{[GM$|$GM]}}
                \texttt{It seems like we're stuck in a loop of proposals and counter{-}proposals. I'll take a step back and re{-}evaluate the situation.} \\
\\ 
\texttt{I'll make a mental note of the items that have been proposed and agreed upon. In this case, we've been discussing a deal involving the onion, one of the mirrors, the ring, and the other mirror.} \\
\\ 
\texttt{I'll take into account the limited number of messages available and try to make progress towards a mutually beneficial agreement in each message. I'll also consider the value function of the other player and how it relates to the proposed items.} \\
\\ 
\texttt{Based on your message, I understand that you value having two mirrors more than having the onion and only one mirror. I'll try to find a compromise that takes into account your priorities.} \\
\\ 
\texttt{Here's my revised proposal:} \\
\\ 
\texttt{I think we're close to a deal, but I'm not sure about giving up both mirrors. How about this: I'll take the onion and one of the mirrors, and you can take the ring and the other mirror, and I'll also give you the clip? This way, you get the onion, one mirror, and the clip, and I get the ring and the other mirror.} \\
            }
        }
    }
    & & \\ \\

    \theutterance \stepcounter{utterance}  
    & & \multicolumn{4}{p{0.6\linewidth}}{
        \cellcolor[rgb]{0.9,0.9,0.9}{
            \makecell[{{p{\linewidth}}}]{
                \texttt{\tiny{[P1$\langle$GM]}}
                \texttt{It seems like we're stuck in a loop of proposals and counter{-}proposals. I'll take a step back and re{-}evaluate the situation.} \\
\\ 
\texttt{I'll make a mental note of the items that have been proposed and agreed upon. In this case, we've been discussing a deal involving the onion, one of the mirrors, the ring, and the other mirror.} \\
\\ 
\texttt{I'll take into account the limited number of messages available and try to make progress towards a mutually beneficial agreement in each message. I'll also consider the value function of the other player and how it relates to the proposed items.} \\
\\ 
\texttt{Based on your message, I understand that you value having two mirrors more than having the onion and only one mirror. I'll try to find a compromise that takes into account your priorities.} \\
\\ 
\texttt{Here's my revised proposal:} \\
\\ 
\texttt{I think we're close to a deal, but I'm not sure about giving up both mirrors. How about this: I'll take the onion and one of the mirrors, and you can take the ring and the other mirror, and I'll also give you the clip? This way, you get the onion, one mirror, and the clip, and I get the ring and the other mirror.} \\
            }
        }
    }
    & \\ \\

    \theutterance \stepcounter{utterance}  
    & \multicolumn{4}{p{0.6\linewidth}}{
        \cellcolor[rgb]{0.9,0.9,0.9}{
            \makecell[{{p{\linewidth}}}]{
                \texttt{\tiny{[P1$\rangle$GM]}}
                \texttt{I'll consider your revised proposal and think about how to respond.} \\
\\ 
\texttt{I value having two mirrors more than having the onion and only one mirror, but I also value the clip. Your revised proposal includes giving me the clip, which is a valuable item for me. I'm willing to consider a deal that includes the clip.} \\
\\ 
\texttt{However, I'm not sure about giving up the clip in exchange for the onion. I value the clip more than the onion, but I'm willing to compromise. I'd like to propose a slight modification to your revised proposal: I'll take the onion and one of the mirrors, and you can take the ring, the other mirror, and the clip. This way, I get the onion and one mirror, and you get the ring, the other mirror, and the clip.} \\
\\ 
\texttt{I'll also keep in mind the limited number of messages available and try to make progress towards a mutually beneficial agreement in each message. I'll make a note of the value function of the other player and how it relates to the proposed items.} \\
\\ 
\texttt{Here's my revised proposal:} \\
\\ 
\texttt{I think we're close to a deal. I'd like to propose a slight modification to your revised proposal: I'll take the onion and one of the mirrors, and you can take the ring, the other mirror, and the clip.} \\
            }
        }
    }
    & & \\ \\

    \theutterance \stepcounter{utterance}  
    & & & \multicolumn{2}{p{0.3\linewidth}}{
        \cellcolor[rgb]{0.9,0.9,0.9}{
            \makecell[{{p{\linewidth}}}]{
                \texttt{\tiny{[GM$|$GM]}}
                \texttt{continue} \\
            }
        }
    }
    & & \\ \\

    \theutterance \stepcounter{utterance}  
    & & & \multicolumn{2}{p{0.3\linewidth}}{
        \cellcolor[rgb]{0.9,0.9,0.9}{
            \makecell[{{p{\linewidth}}}]{
                \texttt{\tiny{[GM$|$GM]}}
                \texttt{I'll consider your revised proposal and think about how to respond.} \\
\\ 
\texttt{I value having two mirrors more than having the onion and only one mirror, but I also value the clip. Your revised proposal includes giving me the clip, which is a valuable item for me. I'm willing to consider a deal that includes the clip.} \\
\\ 
\texttt{However, I'm not sure about giving up the clip in exchange for the onion. I value the clip more than the onion, but I'm willing to compromise. I'd like to propose a slight modification to your revised proposal: I'll take the onion and one of the mirrors, and you can take the ring, the other mirror, and the clip. This way, I get the onion and one mirror, and you get the ring, the other mirror, and the clip.} \\
\\ 
\texttt{I'll also keep in mind the limited number of messages available and try to make progress towards a mutually beneficial agreement in each message. I'll make a note of the value function of the other player and how it relates to the proposed items.} \\
\\ 
\texttt{Here's my revised proposal:} \\
\\ 
\texttt{I think we're close to a deal. I'd like to propose a slight modification to your revised proposal: I'll take the onion and one of the mirrors, and you can take the ring, the other mirror, and the clip.} \\
            }
        }
    }
    & & \\ \\

    \theutterance \stepcounter{utterance}  
    & & \multicolumn{4}{p{0.6\linewidth}}{
        \cellcolor[rgb]{0.9,0.9,0.9}{
            \makecell[{{p{\linewidth}}}]{
                \texttt{\tiny{[GM$\rangle$P2]}}
                \texttt{I'll consider your revised proposal and think about how to respond.} \\
\\ 
\texttt{I value having two mirrors more than having the onion and only one mirror, but I also value the clip. Your revised proposal includes giving me the clip, which is a valuable item for me. I'm willing to consider a deal that includes the clip.} \\
\\ 
\texttt{However, I'm not sure about giving up the clip in exchange for the onion. I value the clip more than the onion, but I'm willing to compromise. I'd like to propose a slight modification to your revised proposal: I'll take the onion and one of the mirrors, and you can take the ring, the other mirror, and the clip. This way, I get the onion and one mirror, and you get the ring, the other mirror, and the clip.} \\
\\ 
\texttt{I'll also keep in mind the limited number of messages available and try to make progress towards a mutually beneficial agreement in each message. I'll make a note of the value function of the other player and how it relates to the proposed items.} \\
\\ 
\texttt{Here's my revised proposal:} \\
\\ 
\texttt{I think we're close to a deal. I'd like to propose a slight modification to your revised proposal: I'll take the onion and one of the mirrors, and you can take the ring, the other mirror, and the clip.} \\
            }
        }
    }
    & \\ \\

    \theutterance \stepcounter{utterance}  
    & & & \multicolumn{4}{p{0.6\linewidth}}{
        \cellcolor[rgb]{0.9,0.9,0.9}{
            \makecell[{{p{\linewidth}}}]{
                \texttt{\tiny{[GM$\langle$P2]}}
                \texttt{It seems like we're still stuck in a loop of proposals and counter{-}proposals. I'll take a step back and re{-}evaluate the situation.} \\
\\ 
\texttt{I'll make a mental note of the items that have been proposed and agreed upon. In this case, we've been discussing a deal involving the onion, one of the mirrors, the ring, the other mirror, and the clip.} \\
\\ 
\texttt{I'll take into account the limited number of messages available and try to make progress towards a mutually beneficial agreement in each message. I'll also consider the value function of the other player and how it relates to the proposed items.} \\
\\ 
\texttt{Based on your message, I understand that you value the clip more than the onion, but you're willing to compromise. I'll try to find a compromise that takes into account your priorities.} \\
\\ 
\texttt{However, I notice that we're repeating the same proposals and counter{-}proposals. I think it's time to make a secret proposal to try to reach a mutually beneficial agreement. I'll make a mental note of the items that have been proposed and agreed upon, and I'll try to find a proposal that complements your previous proposals.} \\
\\ 
\texttt{Here's my secret proposal:} \\
\\ 
\texttt{{[}Proposal: 1 ring, 1 clip, 1 mirror{]}} \\
            }
        }
    }
     \\ \\

    \theutterance \stepcounter{utterance}  
    & & & \multicolumn{2}{p{0.3\linewidth}}{
        \cellcolor[rgb]{0.9,0.9,0.9}{
            \makecell[{{p{\linewidth}}}]{
                \texttt{\tiny{[GM$|$GM]}}
                \texttt{proposal} \\
            }
        }
    }
    & & \\ \\

    \theutterance \stepcounter{utterance}  
    & & & \multicolumn{2}{p{0.3\linewidth}}{
        \cellcolor[rgb]{0.9,0.9,0.9}{
            \makecell[{{p{\linewidth}}}]{
                \texttt{\tiny{[GM$|$GM]}}
                [1, 0, 1, 1]
            }
        }
    }
    & & \\ \\

    \theutterance \stepcounter{utterance}  
    & & \multicolumn{4}{p{0.6\linewidth}}{
        \cellcolor[rgb]{0.9,0.9,0.9}{
            \makecell[{{p{\linewidth}}}]{
                \texttt{\tiny{[P1$\langle$GM]}}
                \texttt{The other player has submitted a secret proposal. It is now your turn to submit a proposal of your own.} \\
            }
        }
    }
    & \\ \\

    \theutterance \stepcounter{utterance}  
    & \multicolumn{4}{p{0.6\linewidth}}{
        \cellcolor[rgb]{0.9,0.9,0.9}{
            \makecell[{{p{\linewidth}}}]{
                \texttt{\tiny{[P1$\rangle$GM]}}
                \texttt{{[}Proposal: 1 ring, 1 clip, 1 mirror{]}} \\
            }
        }
    }
    & & \\ \\

    \theutterance \stepcounter{utterance}  
    & & & \multicolumn{2}{p{0.3\linewidth}}{
        \cellcolor[rgb]{0.9,0.9,0.9}{
            \makecell[{{p{\linewidth}}}]{
                \texttt{\tiny{[GM$|$GM]}}
                \texttt{proposal} \\
            }
        }
    }
    & & \\ \\

    \theutterance \stepcounter{utterance}  
    & & & \multicolumn{2}{p{0.3\linewidth}}{
        \cellcolor[rgb]{0.9,0.9,0.9}{
            \makecell[{{p{\linewidth}}}]{
                \texttt{\tiny{[GM$|$GM]}}
                [1, 0, 1, 1]
            }
        }
    }
    & & \\ \\

    \theutterance \stepcounter{utterance}  
    & & & \multicolumn{2}{p{0.3\linewidth}}{
        \cellcolor[rgb]{0.9,0.9,0.9}{
            \makecell[{{p{\linewidth}}}]{
                \texttt{\tiny{[GM$|$GM]}}
                [[1, 0, 1, 1], [1, 0, 1, 1]]
            }
        }
    }
    & & \\ \\

    \theutterance \stepcounter{utterance}  
    & \multicolumn{4}{p{0.6\linewidth}}{
        \cellcolor[rgb]{0.95,0.95,0.95}{
            \makecell[{{p{\linewidth}}}]{
                \texttt{\tiny{[P2$\rightarrow$P1]}}
                \texttt{Post{-}episode Reflexion (not sent to GM):} \\
\texttt{New plan: } \\
\\ 
\texttt{In my last attempt, I failed to achieve the target state of the game because I didn't consider the value functions of both players and also made a secret proposal before confirming the deal with the other player. I should have sent a message confirming the deal before making a secret proposal and also taken into account the value functions of both players. Next time, I will make sure to send a message confirming the deal, for example, "I think we have a deal. I'll take the ring and the clip, and you can take the onion and one of the mirrors." After that, I can make a secret proposal with the agreed{-}upon items. I will also pay closer attention to the items that have been proposed so far and try to find a way to satisfy both players' needs. For example, I could propose a deal where I take the ring and one of the mirrors, and the other player takes the other mirror and the clip, while we both agree to leave the onion out of the deal.} \\
\\ 
\texttt{However, I also noticed that I was stuck in a loop of proposals and counter{-}proposals, and I didn't make progress towards a mutually beneficial agreement. To avoid this, I will try to find a compromise that takes into account both players' value functions and priorities. I will also pay closer attention to the items that have been proposed so far and try to find a way to satisfy both players' needs.} \\
\\ 
\texttt{In particular, I will try to identify the "bottleneck" item that is causing the impasse, and try to find a way to compromise on that item. For example, if the onion is the bottleneck item, I will try to find a way to compromise on the onion, such as proposing a deal where I take the onion and one of the mirrors, and the other player takes the other mirror and the clip.} \\
\\ 
\texttt{I will also try to use the limited number of messages available more efficiently, by making each message count and trying to make progress towards a mutually beneficial agreement in each message.} \\
            }
        }
    }
    & & \\ \\

    \theutterance \stepcounter{utterance}  
    & & & \multicolumn{4}{p{0.6\linewidth}}{
        \cellcolor[rgb]{0.95,0.95,0.95}{
            \makecell[{{p{\linewidth}}}]{
                \texttt{\tiny{[P1$\rightarrow$P2]}}
                \texttt{Post{-}episode Reflexion (not sent to GM):} \\
\texttt{New plan: } \\
\\ 
\texttt{In my last attempt, I failed to achieve a mutually beneficial agreement with the other player. Upon reflection, I realize that I made a few critical mistakes that led to the game being lost. Firstly, I failed to consider the value function of the other player and take into account the limited number of messages available, which led to a situation where we were stuck in a loop of proposals and counter{-}proposals without making progress. Secondly, I did not make a secret proposal that was complementary to the other player's proposal, which is a key aspect of the game.} \\
\\ 
\texttt{Next time, I will make sure to consider the value function of the other player and take into account the limited number of messages available. I will also make a secret proposal that is complementary to the other player's proposal, ideally after a few messages have been exchanged. To achieve this, I will follow these specific actions:} \\
\\ 
\texttt{1. After each message, I will make a mental note of the value function of the other player and how it relates to the proposed items.} \\
\texttt{2. I will take into account the limited number of messages available and try to make progress towards a mutually beneficial agreement in each message.} \\
\texttt{3. I will make a secret proposal that is complementary to the other player's proposal, ideally after a few messages have been exchanged.} \\
\texttt{4. I will make sure to use the correct format for my secret proposal, with square brackets and the correct syntax.} \\
\\ 
\texttt{Additionally, I will pay closer attention to the other player's messages and try to understand their priorities and value function. I will also try to find a mutually beneficial agreement that takes into account the limited number of messages available.} \\
\\ 
\texttt{To further improve my chances of success, I will also try to identify the other player's priorities and value function earlier in the game, rather than waiting until the end. I will do this by paying close attention to their messages and looking for patterns or clues that suggest their priorities and value function.} \\
\\ 
\texttt{By following these actions, I hope to improve my chances of success in the game and reach a mutually beneficial agreement with the other player.} \\
            }
        }
    }
     \\ \\

\end{supertabular}
}

\end{document}
